\section{A direct factor-filtration route to the open theorem}

The geometric pairing isolates the correct obstruction spaces, but the
auxiliary horizontal unit in the natural pairing is noncanonical.  For the
vanishing theorem there is a more concrete route: calculate the lower
Bockstein directly from the $E_2$ expansion, in the style of
\cite{ARR2026}.

Write $i=np+i'$ and put
\[
 q_1(n,i')=i'^2+(k-1)i'-(n+1)(n+2)\in\mathbb F_p,
 \qquad X=E_2/12\pmod p.
\]
For a $D=0$ cell let
\[
 w_-=k+2i-P,
 \qquad
 \mathcal B_i=\Gamma_{w_-}(\theta^if)\in Q_{k+2i}.
\]
The shifted leading-block candidate is the image in the same quotient of
\begin{equation}\label{eq:shifted-leading-block}
 \mathcal L_i=q_1(n,i')X^p\theta^{i-p}f.
\end{equation}

\begin{proposition}[Finite exact audit]
\label{prop:finite-direct-audit}
In the current Sturm-safe record of $692$ ordinary, nondegenerate
$D=0$--$D=2$ complementary pairs, the following statements hold exactly.
\begin{enumerate}
\item Away from the three affine lines
\begin{equation}\label{eq:three-boundary-lines}
 i'=n+2,
 \qquad i'+n=2p-k-1,
 \qquad i'+n=2p-k,
\end{equation}
one has $\mathcal B_i=\mathcal L_i$ in all $458$ cells.
\item On the three lines in \eqref{eq:three-boundary-lines}, the equality
fails in all $234$ cells, but $\mathcal B_i$ is nonzero in every cell.
\item The $20$ vanishing classes occur off the three lines and are exactly
the cells with $q_1(n,i')=0$.
\item On the line $i'=n+2$, all $78$ tested residuals satisfy
\begin{equation}\label{eq:l3-residual}
 \mathcal B_i-\mathcal L_i
 =c_{p,k,n}X^{2n-2p+k+3}\theta^{p-k}f
 \quad\text{in }Q_{k+2i},
\end{equation}
for an explicitly computed scalar $c_{p,k,n}\in\mathbb F_p^\times$.
\item Put $m=p-2-n$.  On the line $i'+n=2p-k-1$, all $78$
tested residuals lie in the explicit triangular span
\begin{equation}\label{eq:l1-triangular-span}
 \left\langle
 X^a\theta^{2p-k-1-a}f:
 p-k\le a\le p-k+2m+1;
 \ X^p\theta^{p-k-1}f
 \right\rangle .
\end{equation}
On the adjacent line $i'+n=2p-k$, all $78$ tested residuals lie in
\begin{equation}\label{eq:l2-triangular-span}
 \left\langle
 X^a\theta^{2p-k-a}f:
 p-k+1\le a\le p-k+2m+3;
 \ X^p\theta^{p-k}f
 \right\rangle .
\end{equation}
Let $V_m^-$ and $V_m^+$ denote the spans in
\eqref{eq:l1-triangular-span} and \eqref{eq:l2-triangular-span},
respectively, after omitting the displayed $X^p$ endpoint.  In every
tested cell the endpoint class is independent of the corresponding
$V_m^\pm$, and
\begin{align}
 \mathcal B_i-\mathcal L_i
  &\equiv -(m+1)X^p\theta^{p-k-1}f\pmod{V_m^-},
  && i'+n=2p-k-1, \label{eq:l1-endpoint}\\
 \mathcal B_i-\mathcal L_i
  &\equiv -(m+1)X^p\theta^{p-k}f\pmod{V_m^+},
  && i'+n=2p-k. \label{eq:l2-endpoint}
\end{align}
Consequently every one of these $156$ residuals is nonzero.
\end{enumerate}
These statements are finite computations, not assertions for arbitrary
$p$, $k$, and $f$.
\end{proposition}

\begin{proof}
The direct class is obtained by reducing $\theta^if$ modulo $p^2$ by the
echelon basis of $M_{w_-}$, dividing the canonical remainder by $p$, and
then reducing modulo $p$.  The candidate in
\eqref{eq:shifted-leading-block} is reduced by the same basis modulo $p$.
Full Sturm-safe comparison gives $458$ equal and $234$ unequal classes.
The unequal cells are classified by
\eqref{eq:three-boundary-lines} with confusion counts
\[
 \mathrm{TP}=234,\quad \mathrm{FP}=0,
 \quad \mathrm{FN}=0,\quad \mathrm{TN}=458.
\]
The direct remainder vanishes in precisely $20$ cells; each has $q_1=0$.
The boundary residuals are nonzero.  A second full-remainder comparison
tests \eqref{eq:l3-residual} in all $78$ cells on $i'=n+2$; every residual
is proportional to the displayed block and every scalar is nonzero.  This
verifies the proposition over the stated record.  Finally, exact Gaussian
elimination over $\mathbb F_p$ on the complete remainders places every
residual on the other two lines in the spans
\eqref{eq:l1-triangular-span} and \eqref{eq:l2-triangular-span}.  The
restricted echelon solution has support equal to its column rank in every
case.  Omitting the final $X^p$ column lowers the rank by one in all $156$
cases, while the coefficient of that final column is $-(m+1)$ in every
case.  Since $1\le m<p-1$, this is a nonzero scalar and proves the finite
nonvanishing assertions \eqref{eq:l1-endpoint} and
\eqref{eq:l2-endpoint}.
\end{proof}

\begin{theorem}[Explicit periodic first-Hasse-jet formula]
\label{thm:explicit-periodic-first-jet}
Assume that $f\in M_k(\mathbb Z_{(p)})$ is ordinary, that
$4\leq k<p-1$, and that $(n,i')$ is an admissible $D=0$--$D=2$
cross-wedge cell.  Write
\[
 n=p-2-m,
 \qquad
 i'=p-k+m+1+d,
 \qquad
 0\leq d\leq k-2m-1.
\]
Put
\[
 r=p-k+d,
 \qquad L=n+1,
 \qquad \ell_0=p-k+1.
\]
Then $i=np+i'=r+L(p-1)$.

Let $A=E_{p-1}$ and $B=E_{p+1}$ be their standard
$p$-integral Eisenstein lifts, and, for $g\in M_v(\mathbb Z_{(p)})$,
define
\begin{equation}\label{eq:modular-theta-lift}
 \mathcal T_v(g)
 =A\theta g+\frac{v}{12}(B-AE_2)g
 \in M_{v+p+1}(\mathbb Z_{(p)}).
\end{equation}
Let $F_j(f)$ be the modified Serre derivatives, defined by
\[
 F_0(f)=f,
 \qquad F_1(f)=\partial_k f,
 \qquad
 F_{j+1}(f)=\partial_{k+2j}F_j(f)
 -\frac{j(j+k-1)}{144}E_4F_{j-1}(f),
\]
where $\partial_v g=\theta g-(v/12)E_2g$.  At the first ordinary low
point, the Chida--Kaneko congruence gives
\[
 F_{\ell_0}(f)\in M_{k+2\ell_0}(\mathbb Z_{(p)}),
 \qquad
 F_{\ell_0}(f)\equiv\theta^{\ell_0}f\pmod p.
\]
Define a modular lift $G_r$ of $\theta^r f$ as follows.  If
$r<\ell_0$, start with $G_0=f$ and apply \eqref{eq:modular-theta-lift}
$r$ times.  If $r\geq\ell_0$, start with $G_{\ell_0}=F_{\ell_0}(f)$
and apply \eqref{eq:modular-theta-lift} $r-\ell_0$ times.  Its weight is
\[
 v_r=
 \begin{cases}
 k+r(p+1),&r<\ell_0,\\
 k+r(p+1)-\ell_0(p-1),&r\geq\ell_0.
 \end{cases}
\]
If
\[
 a=\frac{w_--v_r}{p-1}
 =\begin{cases}
 k-2m-2,&d=0,\\
 p-2m-d-1,&d\geq1,
 \end{cases}
\]
then $a\geq0$ and
\begin{equation}\label{eq:explicit-lower-weight-lift}
 H_i=A^aG_r\in M_{w_-}(\mathbb Z_{(p)}),
 \qquad H_i\equiv\theta^if\pmod p.
\end{equation}

In $\mathbb F_p[[q]]$, put
\begin{align*}
 \alpha&=\frac{A-1}{p}\pmod p,\\
 \Phi_r(f)&=\frac{\theta^{r+p-1}f-\theta^rf}{p}\pmod p,\\
 \eta_r(f)&=\frac{G_r-\theta^rf}{p}\pmod p.
\end{align*}
These quotients are integral.  The lower Bockstein is the class
\begin{equation}\label{eq:periodic-first-jet-identity}
 \boxed{\quad
 \mathcal B_i
 =\left[
 L\Phi_r(f)-\eta_r(f)-a\alpha\theta^rf
 \right]_{Q_{k+2i}}.
 \quad}
\end{equation}
Thus the minimal-weight projection correction is an explicit first-Hasse-jet
operator; it is not an unspecified choice of characteristic-zero lift.
\end{theorem}

\begin{proof}
The Serre derivative identity rewrites \eqref{eq:modular-theta-lift} as
\[
 \mathcal T_v(g)
 =A\left(\theta g-\frac{v}{12}E_2g\right)
   +\frac{v}{12}Bg,
\]
so it is modular of weight $v+p+1$.  The congruences
$A\equiv1$ and $B\equiv E_2\pmod p$ give
$\mathcal T_v(g)\equiv\theta g\pmod p$.  Together with the
Chida--Kaneko low-point congruence, induction gives
$G_r\equiv\theta^rf\pmod p$ and the displayed formula for $v_r$.
Substitution of $r=p-k+d$ and the two possible filtration branches gives
the stated value of $a$.  The cross-wedge inequalities make both formulas
for $a$ nonnegative.  Hence \eqref{eq:explicit-lower-weight-lift} follows.

For a Fourier index $j$ prime to $p$, write
$j^{p-1}=1+p q_p(j)\pmod {p^2}$.  Since
$i=r+L(p-1)$,
\[
 j^i=j^r(1+pLq_p(j))\pmod {p^2}.
\]
If $p\mid j$, then $r\geq p-k\geq2$, and both sides of the corresponding
divided congruence vanish modulo $p$.  Coefficientwise, therefore,
\begin{equation}\label{eq:fermat-period-jet}
 \frac{\theta^if-\theta^rf}{p}
 \equiv L\Phi_r(f)\pmod p.
\end{equation}
Also $A=1+p\alpha\pmod {p^2}$ and
$G_r=\theta^rf+p\eta_r(f)\pmod {p^2}$, whence
\[
 H_i=A^aG_r
 \equiv\theta^rf+p\bigl(\eta_r(f)+a\alpha\theta^rf\bigr)
 \pmod {p^2}.
\]
Subtracting this identity from \eqref{eq:fermat-period-jet}, dividing by
$p$, and passing to the lower Hasse quotient proves
\eqref{eq:periodic-first-jet-identity}.
\end{proof}

The first Hasse jet in \eqref{eq:periodic-first-jet-identity} is itself
recursive.  Put
\[
 C=\frac{B-AE_2}{p}\pmod p.
\]
Whenever $G_{t+1}=\mathcal T_{v_t}(G_t)$, direct expansion modulo $p^2$
gives
\begin{equation}\label{eq:eta-recursion}
 \eta_{t+1}
 =\theta\eta_t+\alpha\theta^{t+1}f
   +\frac{v_t}{12}C\theta^tf.
\end{equation}
The initial jet is $\eta_0=0$ on the first branch and
$(F_{\ell_0}(f)-\theta^{\ell_0}f)/p$ on the second.  Equations
\eqref{eq:periodic-first-jet-identity} and \eqref{eq:eta-recursion}
therefore reduce the remaining block identities to a finite symbolic
calculation in three explicit divided $q$-series.

\begin{lemma}[Unit-block endpoint coefficient]
\label{lem:unit-block-endpoint}
Suppose that $i=np+i'$ lies on either of the two triangular lines in
\eqref{eq:three-boundary-lines}, and put $m=p-2-n$.  Define
\[
 \kappa_A
 =\frac1p\binom{i}{i-p}
   \prod_{r=i-p+k}^{i+k-1}r\pmod p.
\]
Then
\[
 \kappa_A=-n(n+1)=-(m+1)(m+2)\pmod p,
\]
and the coefficient of the $X^p\theta^{i-p}f$ unit block after subtracting
the shifted coefficient $q_1(n,i')$ is
\begin{equation}\label{eq:unit-minus-shifted}
 \kappa_A+i(i+k-1)-q_1(n,i')=-2(m+1)\pmod p.
\end{equation}
\end{lemma}

\begin{proof}
Lucas's theorem gives
$\binom{i}{i-p}=\binom{i}{p}\equiv n\pmod p$.  On the first triangular
line one has $i'=p-k+m+1$, and on the second one has
$i'=p-k+m+2$.  In either case the product defining $\kappa_A$ consists
of $p$ consecutive integers and its unique multiple of $p$ is
$(n+1)p$.  After division by $p$, the remaining factors run through
$\mathbb F_p^\times$.  Wilson's theorem therefore gives
\[
 \frac1p\prod_{r=i-p+k}^{i+k-1}r
 \equiv -(n+1)\pmod p,
\]
which proves the formula for $\kappa_A$.  Since
$i(i+k-1)\equiv i'(i'+k-1)$ and
\[
 q_1(n,i')=i'(i'+k-1)-(n+1)(n+2),
\]
their difference is
\[
 -n(n+1)+(n+1)(n+2)=2(n+1)=-2(m+1)\pmod p.
\]
\end{proof}

\begin{proposition}[Uniform separation of the triangular endpoint]
\label{prop:triangular-endpoint-independent}
Assume that $f$ is ordinary at $p$, that $4\leq k<p-1$, and that a
$D=0$--$D=2$ cross-wedge cell lies on one of the two triangular lines in
\eqref{eq:three-boundary-lines}.  Put $m=p-2-n$.  Then
\[
 1\leq m\leq \frac{k-4}{2}.
\]
On the line $i'+n=2p-k-1$, the class
\[
 X^p\theta^{p-k-1}f
\]
is not in $V_m^-+M_{w_-}(\mathbb F_p)$.  On the line
$i'+n=2p-k$, the class
\[
 X^p\theta^{p-k}f
\]
is not in $V_m^++M_{w_-}(\mathbb F_p)$.  In particular, the endpoint
classes in \eqref{eq:l1-endpoint} and \eqref{eq:l2-endpoint} are
independent of the preceding triangular spans for every admissible
$p,k,m$, not only in the finite audit.
\end{proposition}

\begin{proof}
Let $\ell_0=p-k+1$ be the first low point of the ordinary mod-$p$ theta
cycle.  For $0\leq t<p$, its filtration is
\[
 \omega_p(\theta^t f)=
 \begin{cases}
  k+t(p+1),&t<\ell_0,\\
  k+t(p+1)-(p-k+1)(p-1),&t\geq\ell_0.
 \end{cases}
\]
Also
\[
 \omega_p(X^a\theta^t f)=\omega_p(\theta^t f)+a(p+1).
\]

On the first triangular line, every generator of $V_m^-$ has
$a+t=2p-k-1$ and
\[
 p-2m-2\leq t\leq p-1.
\]
The cross-wedge inequalities give $2m+2\leq k-2$, hence
$t\geq p-k+2>\ell_0$.  All these generators therefore have exact
filtration
\[
 F_-^{\rm old}=k+(2p-k-1)(p+1)-(p-k+1)(p-1)
                 =p^2+p-k.
\]
The endpoint has $a=p$ and $t=p-k-1<\ell_0$, so its exact filtration is
\[
 F_-^{\rm end}=k+(2p-k-1)(p+1)=p(2p+1-k)-1.
\]
Thus
\[
 F_-^{\rm end}-F_-^{\rm old}=(p-k+1)(p-1)>0.
\]
Moreover a direct substitution in
$w_-=k+2i-p(p-1)$ gives
\[
 F_-^{\rm old}-w_-=2(m+1)(p-1)>0.
\]

On the second triangular line, every generator of $V_m^+$ has
$a+t=2p-k$ and
\[
 p-2m-3\leq t\leq p-1.
\]
Here $2m+3\leq k-1$, so again $t\geq\ell_0$.  Their common exact
filtration is
\[
 F_+^{\rm old}=k+(2p-k)(p+1)-(p-k+1)(p-1)
                 =p^2+2p-k+1.
\]
The endpoint has $a=p$ and $t=p-k<\ell_0$, and hence
\[
 F_+^{\rm end}=k+(2p-k)(p+1)=p(2p+2-k).
\]
Again
\[
 F_+^{\rm end}-F_+^{\rm old}=(p-k+1)(p-1)>0,
 \qquad
 F_+^{\rm old}-w_-=(2m+3)(p-1)>0.
\]

If an endpoint belonged to the sum of the preceding span and
$M_{w_-}$, its filtration would be at most the larger of $w_-$ and the
filtrations of the preceding generators.  The strict inequalities above
contradict the exact endpoint filtration.  This proves both assertions.
\end{proof}

The exact audit further separates the endpoint in
\eqref{eq:l1-endpoint}--\eqref{eq:l2-endpoint}.  The triangular
$B$-block sum has zero distinguished-endpoint coefficient, while the
explicit first-Hasse-jet correction in
\eqref{eq:periodic-first-jet-identity} has coefficient $m+1$ in all
$156$ tested cells.  Together with \eqref{eq:unit-minus-shifted}, this gives
\[
 -2(m+1)+0+(m+1)=-(m+1).
\]
Thus the open uniform statement on these two lines is now the
associated-graded simplification
\begin{equation}\label{eq:projection-endpoint-target}
 \Pi_i\equiv (m+1)X^p\theta^{p-k-\varepsilon}f
 \pmod{V_m^\varepsilon},
 \qquad \varepsilon=
 \begin{cases}
  1,&i'+n=2p-k-1,\\
  0,&i'+n=2p-k,
 \end{cases}
\end{equation}
where $V_m^1=V_m^-$ and $V_m^0=V_m^+$.  The independence of the
displayed endpoint is Proposition~\ref{prop:triangular-endpoint-independent};
the first-Hasse-jet formula itself is Theorem~\ref{thm:explicit-periodic-first-jet}.
Only its simplification to the displayed coefficient remains open.

There is a parallel uniform statement on the third line.  It does not
determine the scalar in \eqref{eq:l3-residual}, but it proves that the
displayed block is a genuine, nonzero quotient direction.

\begin{proposition}[The third-line block is nonzero]
\label{prop:l3-block-nonzero}
Assume the hypotheses of Proposition
\ref{prop:triangular-endpoint-independent} and suppose that $i'=n+2$.
Put $m=p-2-n$.  Then $i'=p-m$, $1\leq m\leq(k-4)/2$, and
\[
 Y_{p,k,m}:=X^{k-2m-1}\theta^{p-k}f
\]
has exact filtration
\[
 \omega_p(Y_{p,k,m})=w_-+(p-1).
\]
In particular, its image in the lower Hasse quotient is nonzero.  Also
\[
 q_1(n,i')=-km\not\equiv0\pmod p.
\]
\end{proposition}

\begin{proof}
The equality $i'=n+2$ and the definition $m=p-2-n$ give $i'=p-m$.
The cross-wedge inequalities give $1\leq m\leq(k-4)/2$.  Since
$p-k$ lies strictly before the first low point $p-k+1$ of the ordinary
theta cycle,
\[
 \omega_p(\theta^{p-k}f)=k+(p-k)(p+1).
\]
Multiplication by $X^{k-2m-1}$ raises exact filtration by
$(k-2m-1)(p+1)$.  Hence
\[
 \omega_p(Y_{p,k,m})
 =k+(p-2m-1)(p+1).
\]
On this line
\[
 i=(p-2-m)p+(p-m)=p^2-p-m(p+1),
\]
so
\[
 w_-=k+2i-p(p-1)=k+p^2-p-2m(p+1).
\]
Subtracting gives $\omega_p(Y_{p,k,m})-w_-=p-1$.  Thus
$Y_{p,k,m}$ cannot lie in $M_{w_-}(\mathbb F_p)$ and its quotient class
is nonzero.  Finally, direct substitution of $n=p-2-m$ and $i'=p-m$ in
$q_1$ gives $q_1=-km$.  Both $k$ and $m$ lie strictly between $0$ and
$p$, so this scalar is nonzero.
\end{proof}

The scalar in \eqref{eq:l3-residual}, when that identity holds, is
therefore uniquely determined.  In canonical quotient coordinates it is
\begin{equation}\label{eq:l3-scalar-definition}
 c_{p,k,m}(f)=
 \frac{\lambda(\mathcal B_i-\mathcal L_i)}
      {\lambda(Y_{p,k,m})},
\end{equation}
where $\lambda$ is any linear functional nonzero on the one-dimensional
line spanned by $Y_{p,k,m}$.  The value is independent of the choice of
$\lambda$.  Exact reconstruction of the computed values identifies the
following closed candidate:
\begin{equation}\label{eq:l3-closed-scalar}
 c^*_{p,k,m}=-m\,(k-2m-2)!\pmod p.
\end{equation}

\begin{proposition}[The closed third-line scalar is a unit]
\label{prop:l3-closed-scalar-unit}
For every admissible $p,k,m$, the value in
\eqref{eq:l3-closed-scalar} belongs to $\mathbb F_p^\times$.
\end{proposition}

\begin{proof}
The cross-wedge inequalities give
\[
 1\leq m\leq\frac{k-4}{2},
 \qquad
 2\leq k-2m-2\leq k-4<p.
\]
Thus neither $m$ nor any factor in $(k-2m-2)!$ is divisible by $p$.
Their product is nonzero modulo $p$.
\end{proof}

The finite evidence for this stronger statement is broader than the
discovery table.  A blind extension through $p=61$ checked $464$
third-line cells and found no zero scalar, no failure of the one-block
shape.  A separate weight-$24$ test, where the cusp-form space is
two-dimensional, checked $1{,}860$ cells from $186$ ordinary forms and
linear combinations through $p=43$.  In both audits every computed scalar
was exactly \eqref{eq:l3-closed-scalar}; there were no mismatches.  The
scalar was also independent of the form for fixed $p,k,m$.  These exact
computations identify the formula and strongly test it; they do not replace
the symbolic block reduction below.

The generic identity and the third-line formula now have one common target.
For the parameters of Theorem~\ref{thm:explicit-periodic-first-jet}, put
\[
 Z_d=X^d\theta^{p-k}f.
\]
A direct filtration calculation gives
\begin{equation}\label{eq:zd-filtration-gap}
 \omega_p(Z_d)-w_-
 =(d-k+2m+2)(p-1).
\end{equation}
Consequently $Z_d\in M_{w_-}(\mathbb F_p)$ for
$d\leq k-2m-2$, whereas $Z_{k-2m-1}=Y_{p,k,m}$ is the nonzero direction
of Proposition~\ref{prop:l3-block-nonzero}.  It is therefore enough to
prove the following single symbolic reduction.

\begin{conjecture}[Factorial block reduction]
\label{conj:factorial-block-reduction}
For $2\leq d\leq k-2m-1$, the explicit first-Hasse-jet class satisfies
\begin{equation}\label{eq:factorial-block-reduction}
 \mathcal B_i-\mathcal L_i
 =-m(d-1)![Z_d]
 \quad\text{in }Q_{k+2i}.
\end{equation}
\end{conjecture}

For $2\leq d\leq k-2m-2$, equation
\eqref{eq:zd-filtration-gap} makes the right side of
\eqref{eq:factorial-block-reduction} zero, proving the generic shifted
identity.  At $d=k-2m-1$, the same equation becomes
\eqref{eq:l3-residual} with the explicit scalar
\eqref{eq:l3-closed-scalar}, which is a unit by Proposition
\ref{prop:l3-closed-scalar-unit}.  Thus one factorial reduction proves both
the generic identity and the complete third-line assertion.

Theorem~\ref{thm:explicit-periodic-first-jet} replaces the formerly
unspecified minimal-weight projection compatibility square by an exact operator
formula.  What remains is the formal reduction of
\eqref{eq:periodic-first-jet-identity} to
\eqref{eq:factorial-block-reduction}, together with the two triangular
endpoint coefficients.  The audit identifies those associated-graded
coefficients as the unit $-(m+1)$.  Proposition
\ref{prop:triangular-endpoint-independent} already supplies endpoint
independence; it must not be counted among the remaining conjectural inputs.

